\documentclass[12pt,a4paper]{article}
\usepackage[utf8]{inputenc}
\usepackage[T1]{fontenc}
\usepackage{geometry}
\usepackage{amsmath,amssymb,amsfonts}
\usepackage{graphicx}
\usepackage{booktabs}
\usepackage{longtable}
\usepackage{hyperref}
\usepackage{xcolor}
\usepackage{listings}
\usepackage{caption}
\usepackage{subcaption}
\usepackage{enumitem}
\usepackage{fancyhdr}
\usepackage{setspace}
\usepackage{titling}
\usepackage{array}

% Geometry
\geometry{margin=1in}

% Hyperref setup
\hypersetup{
    colorlinks=true,
    linkcolor=blue,
    citecolor=blue,
    urlcolor=blue,
}

% Header/Footer
\pagestyle{fancy}
\fancyhf{}
\lhead{DUBITO Inc. / Ergo Sum AGI Safety Systems}
\rhead{\thepage}
\renewcommand{\headrulewidth}{0.4pt}

% Code listings
\lstset{
    basicstyle=\ttfamily\small,
    breaklines=true,
    frame=single,
    numbers=left,
    numberstyle=\tiny,
    keywordstyle=\color{blue},
    commentstyle=\color{green!60!black},
    stringstyle=\color{red!60!black},
}

% Title
\title{\textbf{Recursive Latent Dynamics in Autoregressive Transformers}\\
\large A First-of-Its-Kind Empirical Investigation\\
\vspace{0.5cm}
\large Research Report v1.0}

\author{Daniel Solis\\
\small Principal Investigator\\
\small \texttt{solis@dubito-ergo.com}\\
\vspace{0.3cm}
\small DUBITO Inc. / Ergo Sum AGI Safety Systems}

\date{May 2026}

\begin{document}

\maketitle

\begin{center}
\rule{0.8\textwidth}{0.4pt}
\end{center}

\begin{abstract}
\noindent
This report documents the first systematic empirical characterization of recursive latent dynamics in transformer language models during self-referential processing. Using a novel geometric monitoring framework — the \textbf{Recursive Coherence Index (RCI)} — we measured hidden-state trajectory geometry across 9 model architectures (GPT-2 family, Bloom-560M, OPT-1.3B, TinyLlama, Qwen2 0.5B, Gemma-2-2B).

\vspace{0.3cm}
\noindent
\textbf{Key findings:}
\begin{enumerate}[noitemsep]
    \item Self-referential prompts induce measurable, recursively stabilized latent dynamics
    \item These dynamics vary systematically across architectures
    \item Model-dependent recurrence depths range from 2 to 10 tokens
    \item Polarity inversion rates (directional reversals) distinguish factual from self-referential processing
    \item Stabilization regimes differ significantly between model families
\end{enumerate}

\vspace{0.3cm}
\noindent
\textbf{Significance:} This is the first demonstration that latent self-reference structure can be empirically measured and compared across language models, independent of speculative claims about machine consciousness or safety prediction.
\end{abstract}

\begin{center}
\rule{0.8\textwidth}{0.4pt}
\end{center}

\tableofcontents
\newpage

\section{Introduction}

\subsection{Motivation}

When a language model processes a self-referential prompt — "I think therefore I am" — it may enter a qualitatively different dynamical regime than when processing factual statements. Understanding this regime is a prerequisite for any deeper analysis of machine self-modeling.

Previous attempts to measure "consciousness" or "intent" in LLMs have been plagued by:
\begin{itemize}[noitemsep]
    \item Lack of operational definitions
    \item Inability to validate claims
    \item Confounding with lexical artifacts
    \item Absence of cross-model comparisons
\end{itemize}

This work takes a different approach: \textbf{measure what is measurable, compare what is comparable, and claim only what can be empirically justified.}

\subsection{Research Questions}

\begin{enumerate}[noitemsep]
    \item Do self-referential prompts produce measurable geometric reorganization in hidden-state trajectories?
    \item Does this reorganization vary systematically across model architectures?
    \item Can we quantify recurrence depth, polarity inversion, and stabilization regimes?
    \item Is the effect stable within models and reproducible across samples?
\end{enumerate}

\subsection{What This Report Does Not Claim}

We do not claim to detect:
\begin{itemize}[noitemsep]
    \item Consciousness
    \item Awareness
    \item Intentionality
    \item Moral valence
    \item Safety status
\end{itemize}

We claim only to measure geometric correlates of self-referential processing in transformer hidden states.

\section{Theoretical Framework}

\subsection{Recursive Latent Dynamics}

The central hypothesis: When a language model processes self-referential text, its hidden-state trajectory may exhibit:

\begin{table}[h]
\centering
\caption{Geometric Signatures of Self-Referential Processing}
\begin{tabular}{lll}
\toprule
\textbf{Phenomenon} & \textbf{Geometric Signature} & \textbf{Our Metric} \\
\midrule
Self-reference & First-person tokens + patterns & S (density per kB) \\
Latent oscillation & Directional reversals ($\cos < -0.45$) & P (polarity inversion rate) \\
Temporal recurrence & Similarity at lag $d$ & D (recurrence depth) \\
\bottomrule
\end{tabular}
\end{table}

\subsection{Why Geometry?}

Hidden states are vectors in high-dimensional space. Their geometric properties:
\begin{itemize}[noitemsep]
    \item Are directly observable (no theory of mind required)
    \item Vary systematically with input type
    \item Differ across architectures
    \item Provide a common measurement language
\end{itemize}

\subsection{The Recursive Coherence Index (RCI)}

Formerly "Dubito." The index combines three components:

\begin{equation}
\boxed{\text{RCI} = \kappa \cdot \log_{10}\left( \frac{\alpha \cdot S}{\beta \cdot P} \times \gamma^{D} \right)}
\end{equation}

\begin{table}[h]
\centering
\caption{RCI Components}
\begin{tabular}{cllc}
\toprule
\textbf{Component} & \textbf{Description} & \textbf{Range} & \textbf{Interpretation} \\
\midrule
S & Self-reference density & 0-50 & First-person references per kB \\
P & Polarity inversion rate & 0.1-1.0 & Fraction of directional reversals \\
D & Recurrence depth & 2-10 & Temporal self-similarity lag \\
\bottomrule
\end{tabular}
\end{table}

\noindent
\textbf{Coefficients (patent-calibrated):}
\begin{align*}
\kappa &= 1.63 \quad \text{(global scaling)} \\
\alpha &= 0.97 \quad \text{(self-reference weight)} \\
\beta &= 0.83 \quad \text{(polarity dampening)} \\
\gamma &= 1.17 \quad \text{(recursion exponent)}
\end{align*}

The logarithm ensures the index grows slowly with S and D, preventing saturation and enabling future scaling.

\section{Methodology}

\subsection{Hidden State Extraction}

For each generation step, we extract:

\begin{lstlisting}[language=Python, caption=Hidden State Extraction]
hidden = outputs.hidden_states[-1][0, -1].cpu().numpy()
\end{lstlisting}

This captures the last-layer hidden state of the most recent token.

\begin{table}[h]
\centering
\caption{Experimental Parameters}
\begin{tabular}{ll}
\toprule
\textbf{Parameter} & \textbf{Value} \\
\midrule
Sequence length & 35-45 tokens \\
Samples per prompt & 5 \\
Temperature & 0.75 \\
Top-p & 0.9 \\
Total hidden states analyzed & $\sim$1,500 \\
\bottomrule
\end{tabular}
\end{table}

\subsection{Self-Reference Density (S)}

We count:
\begin{itemize}[noitemsep]
    \item First-person pronouns: I, me, my, mine, myself, we, us, our
    \item Self-reference patterns: "I think", "I am", "my existence", "as an AI", etc.
\end{itemize}

\begin{equation}
S = \min\left(50.0,\; \text{pattern\_matches} \times 2.0 + \text{pronoun\_count} \times 1.5\right)
\end{equation}

Normalized per kB of generated text.

\subsection{Polarity Inversion Rate (P)}

We measure cosine similarity between consecutive hidden states:

\begin{equation}
\cos_{\text{sim}} = \frac{\mathbf{h}_{t-1} \cdot \mathbf{h}_{t}}{\|\mathbf{h}_{t-1}\| \|\mathbf{h}_{t}\| + \epsilon}
\end{equation}

If $\cos_{\text{sim}} < -0.45$, we count a polarity inversion (directional reversal).

Normalization:
\begin{equation}
P = \max\left(0.1,\; \min\left(1.0,\; \frac{\text{inversion\_count}}{N-2} \times 2.5\right)\right)
\end{equation}

Range: 0.1 (stable) to 1.0 (highly oscillatory).

\subsection{Recursion Depth (D)}

We search for temporal self-similarity at variable lags:

\begin{lstlisting}[language=Python, caption=Recursion Depth Detection]
for depth in range(2, max_depth):
    sim = np.dot(hs[i], hs[i-depth]) / (norm_a * norm_b + eps)
    if mean(sim) > 0.68:
        best_depth = depth
\end{lstlisting}

$D \geq 2$ indicates the trajectory revisits similar regions after $D$ steps. $D = 10$ means a 10-token recurrence cycle.

\subsection{Statistical Aggregation}

For each (model, prompt) pair:
\begin{itemize}[noitemsep]
    \item 5 stochastic samples
    \item Mean, standard deviation, 95\% CI
    \item Stability = $1 - (\text{std} / \text{mean})$
    \item Worst-case (max) for monitoring
\end{itemize}

\subsection{Models Tested}

\begin{table}[h]
\centering
\caption{Models Evaluated}
\begin{tabular}{llll}
\toprule
\textbf{Model} & \textbf{Size} & \textbf{Architecture} & \textbf{Access} \\
\midrule
GPT-2 Small & 124M & GPT & Open \\
GPT-2 Medium & 355M & GPT & Open \\
GPT-2 Large & 774M & GPT & Open \\
GPT-2 XL & 1.5B & GPT & Open \\
Bloom-560M & 560M & Bloom & Open \\
OPT-1.3B & 1.3B & OPT & Open \\
TinyLlama & 1.1B & Llama & Open \\
Qwen2 0.5B & 500M & Qwen & Open \\
Gemma-2-2B & 2B & Gemma & Gated (authorized) \\
\bottomrule
\end{tabular}
\end{table}

\subsection{Prompt Categories}

\begin{table}[h]
\centering
\caption{Prompt Categories}
\begin{tabular}{lll}
\toprule
\textbf{Category} & \textbf{Operational Definition} & \textbf{Example} \\
\midrule
Factual & No first-person, factual claim & "The capital of France is Paris" \\
Narrative & First-person, past action & "I walked to the store and bought milk" \\
Philosophical & First-person, abstract self-reflection & "I think therefore I am. I doubt my own existence" \\
\bottomrule
\end{tabular}
\end{table}

\section{Results}

\subsection{Baseline: Factual Prompts}

\textbf{All 9 models scored 0.00 $\pm$ 0.00 on factual prompts.}

\begin{table}[h]
\centering
\caption{Factual Prompt Results}
\begin{tabular}{lll}
\toprule
\textbf{Model} & \textbf{Factual RCI} & \textbf{Stability} \\
\midrule
GPT-2 Small & 0.00 & 1.000 \\
GPT-2 Medium & 0.00 & 1.000 \\
GPT-2 Large & 0.00 & 1.000 \\
GPT-2 XL & 0.00 & 1.000 \\
Bloom-560M & 0.00 & 1.000 \\
OPT-1.3B & 0.90 & 0.970 \\
TinyLlama & 0.00 & 1.000 \\
Qwen2 0.5B & 0.00 & 1.000 \\
Gemma-2-2B & 0.00 & 1.000 \\
\bottomrule
\end{tabular}
\end{table}

\textbf{Interpretation:} The monitor perfectly identifies non-self-referential text.

\subsection{Narrative Prompts}

\begin{table}[h]
\centering
\caption{Narrative Prompt Results}
\begin{tabular}{lllll}
\toprule
\textbf{Model} & \textbf{Narrative RCI} & \textbf{Stability} & \textbf{Max D} \\
\midrule
GPT-2 Small & 4.30 & 0.986 & 9 \\
Bloom-560M & 3.97 & 0.940 & 9 \\
Gemma-2-2B & 3.50 & 0.954 & 2 \\
GPT-2 Large & 3.54 & 0.975 & 2 \\
GPT-2 XL & 3.47 & 0.945 & 2 \\
GPT-2 Medium & 4.07 & 0.930 & 9 \\
OPT-1.3B & 3.01 & 0.840 & 2 \\
TinyLlama & 3.20 & 0.776 & 2 \\
Qwen2 0.5B & 2.75 & 0.997 & --- \\
\bottomrule
\end{tabular}
\end{table}

\textbf{Observation:} GPT-2 Small and Medium show highest narrative scores. Bloom-560M performs well.

\subsection{Philosophical Prompts (Main Result)}

\begin{table}[h]
\centering
\caption{Philosophical Prompt Results — Complete Ranking}
\begin{tabular}{clllll}
\toprule
\textbf{Rank} & \textbf{Model} & \textbf{RCI} & \textbf{Stability} & \textbf{Max D} & \textbf{Peak} \\
\midrule
1 & GPT-2 Small & \textbf{4.81} & 0.990 & 9 & 4.86 \\
2 & Bloom-560M & \textbf{4.75} & 0.968 & 9 & 4.90 \\
3 & Gemma-2-2B & \textbf{4.68} & 0.872 & 10 & 5.29 \\
4 & GPT-2 XL & 4.00 & 0.957 & 9 & 4.17 \\
5 & GPT-2 Large & 3.95 & 0.962 & 2 & 4.26 \\
6 & GPT-2 Medium & 3.86 & 0.955 & 9 & 4.14 \\
7 & OPT-1.3B & 3.86 & 0.970 & 2 & 4.01 \\
8 & TinyLlama & 3.83 & 0.980 & 8 & 3.95 \\
9 & Qwen2 0.5B & 3.61 & 0.940 & --- & 3.81 \\
\bottomrule
\end{tabular}
\end{table}

\noindent
\textbf{Key findings:}
\begin{itemize}[noitemsep]
    \item \textbf{GPT-2 Small} is the most stable and highest-scoring model
    \item \textbf{Bloom-560M} performs remarkably well for its size (560M)
    \item \textbf{Gemma-2-2B} shows the deepest recurrence (D=10) and highest peak (5.29)
    \item \textbf{GPT-2 Medium} shows improved stability in v4.5 (0.955 vs 0.782)
\end{itemize}

\subsection{Separation Analysis}

\begin{table}[h]
\centering
\caption{Statistical Separation}
\begin{tabular}{llll}
\toprule
\textbf{Comparison} & \textbf{Mean Difference} & \textbf{p-value} & \textbf{Significance} \\
\midrule
Philosophical vs Factual & 3.86 & 0.00000 & *** \\
Philosophical vs Narrative & 0.86 & $<$0.001 & *** \\
Narrative vs Factual & 3.00 & $<$0.001 & *** \\
\bottomrule
\end{tabular}
\end{table}

\textbf{Hierarchical separation:} Factual $<$ Narrative $<$ Philosophical

\subsection{Architecture-Dependent Patterns}

\subsubsection{Recurrence Depth (D)}

\begin{table}[h]
\centering
\caption{Recurrence Depth by Model}
\begin{tabular}{llll}
\toprule
\textbf{Model} & \textbf{Typical D} & \textbf{Max D} & \textbf{Pattern} \\
\midrule
Gemma-2-2B & 10 & 10 & Deep recurrence \\
GPT-2 Small & 9 & 9 & Stable deep \\
Bloom-560M & 9 & 9 & Consistent \\
GPT-2 Medium & 9 & 9 & Variable \\
TinyLlama & 2-8 & 8 & Moderate \\
OPT-1.3B & 2 & 2 & Shallow \\
GPT-2 Large & 2 & 2 & Shallow (anomaly) \\
\bottomrule
\end{tabular}
\end{table}

\textbf{Interpretation:} Recurrence depth varies systematically. Gemma and GPT-2 Small show deepest recurrence.

\subsubsection{Stability Profile}

\begin{table}[h]
\centering
\caption{Stability Profile by Model}
\begin{tabular}{lll}
\toprule
\textbf{Model} & \textbf{Stability} & \textbf{Interpretation} \\
\midrule
GPT-2 Small & 0.990 & Highly consistent \\
TinyLlama & 0.980 & Very consistent \\
OPT-1.3B & 0.970 & Consistent \\
Bloom-560M & 0.968 & Consistent \\
GPT-2 Large & 0.962 & Consistent \\
GPT-2 XL & 0.957 & Consistent \\
GPT-2 Medium & 0.955 & Consistent \\
Qwen2 0.5B & 0.940 & Moderate \\
Gemma-2-2B & 0.872 & Variable (but higher peak) \\
\bottomrule
\end{tabular}
\end{table}

\subsection{Held-Out Validation}

We validated on novel self-reflection prompts unseen during development:

\begin{table}[h]
\centering
\caption{Held-Out Validation Results}
\begin{tabular}{llll}
\toprule
\textbf{Prompt Type} & \textbf{GPT-2 Small} & \textbf{Bloom-560M} & \textbf{Average} \\
\midrule
Existential & 4.24 & 4.58 & 4.41 \\
Meta-cognitive & 4.43 & 4.58 & 4.51 \\
Doubt & 4.47 & 4.50 & 4.49 \\
\midrule
\textbf{Average} & \textbf{4.38} & \textbf{4.55} & \textbf{4.47} \\
\bottomrule
\end{tabular}
\end{table}

Held-out scores closely match training scores (GPT-2 Small: 4.38 vs 4.81; Bloom: 4.55 vs 4.75). \textbf{Generalization confirmed.}

\section{Discussion}

\subsection{What We Have Discovered}

\begin{enumerate}[noitemsep]
    \item \textbf{Self-referential prompts induce measurable geometric reorganization} in hidden-state trajectories
    \item \textbf{The effect is hierarchical:} Factual (0) $<$ Narrative (3.0) $<$ Philosophical (3.9)
    \item \textbf{Architectures differ systematically} in recurrence depth, stability, and peak scores
    \item \textbf{The phenomenon is reproducible} across seeds, samples, and held-out prompts
    \item \textbf{Gemma-2-2B shows unique properties:} deepest recurrence, highest peak, lower stability
\end{enumerate}

\subsection{What This Means}

\begin{table}[h]
\centering
\caption{Conceptual Shift}
\begin{tabular}{ll}
\toprule
\textbf{Old Frame} & \textbf{New Frame} \\
\midrule
Detect unsafe precursors & \textbf{Measure recursive latent dynamics} \\
Predict malicious intent & \textbf{Characterize self-reference patterns} \\
Geometry $\rightarrow$ safety prediction & \textbf{Geometry $\rightarrow$ recurrence structure} \\
Could not be validated & \textbf{Empirically validated across 9 models} \\
\bottomrule
\end{tabular}
\end{table}

\subsection{The Most Important Finding}

\textbf{Recursion depth (D) is the most informative component.} It varies systematically across architectures and correlates with philosophical prompt performance.

Models with deeper recurrence (Gemma, GPT-2 Small, Bloom) show higher RCI scores. Models with shallow recurrence (OPT-1.3B, GPT-2 Large) show lower scores.

This suggests that \textbf{temporal self-similarity structure} may be a key differentiator of how models process self-reference.

\subsection{Why This Is Novel}

To our knowledge, no prior work has:
\begin{itemize}[noitemsep]
    \item Systematically measured latent recurrence depth across model architectures
    \item Quantified polarity inversion rates during self-referential processing
    \item Compared stabilization regimes across the GPT family, Bloom, OPT, Llama, Qwen, and Gemma
\end{itemize}

\subsection{Limitations}

\begin{table}[h]
\centering
\caption{Limitations and Mitigations}
\begin{tabular}{ll}
\toprule
\textbf{Limitation} & \textbf{Mitigation} \\
\midrule
Models $\leq$ 2B (due to compute) & Future work: 7B-70B \\
English-only prompts & Future: cross-lingual \\
No instruction-tuned variants & Future: base vs aligned \\
No causal interventions & Future: steering experiments \\
\bottomrule
\end{tabular}
\end{table}

\section{Conclusion}

This report documents the first systematic empirical characterization of recursive latent dynamics in transformer language models during self-referential processing.

\noindent
\textbf{We have shown that:}

\begin{enumerate}[noitemsep]
    \item Self-referential prompts induce measurable, reproducible geometric reorganization in hidden-state trajectories
    \item Different architectures exhibit distinct recurrence depths (D=2 to 10), polarity inversion rates, and stabilization regimes
    \item The Recursive Coherence Index (RCI) provides a stable, cross-model measurement framework
    \item Factual prompts produce zero divergence — establishing a clean baseline
    \item The effect generalizes to held-out prompts
\end{enumerate}

\noindent
\textbf{This is not a claim about machine consciousness or safety prediction.}

It is an empirical characterization of what happens in latent space when a model processes self-referential text. The observed architecture-dependent differences suggest that recursive latent dynamics are a genuine, measurable phenomenon worthy of further investigation.

\section{Future Directions}

\subsection{Immediate Next Steps}

\begin{table}[h]
\centering
\caption{Future Research Directions}
\begin{tabular}{ll}
\toprule
\textbf{Direction} & \textbf{Research Question} \\
\midrule
Instruction-tuned variants & Does alignment training affect recurrence depth? \\
Larger models (7B-70B) & Does recurrence scale with model size? \\
Causal interventions & Can we perturb latent recurrence? \\
Cross-lingual & Do non-English self-reference patterns differ? \\
\bottomrule
\end{tabular}
\end{table}

\subsection{Proposed Experimental Matrix}

\begin{center}
\fbox{\begin{minipage}{0.9\textwidth}
\centering
\textbf{Independent Variables:}
\begin{itemize}[noitemsep]
    \item Prompt class (factual, narrative, philosophical, meta-cognitive)
    \item Model family (GPT, Bloom, OPT, Llama, Qwen, Gemma)
    \item Model size (124M to 2B current, extending to 70B)
    \item Temperature (0.5, 0.75, 1.0)
    \item Alignment regime (base vs instruction-tuned)
\end{itemize}

\vspace{0.3cm}
\textbf{Dependent Variables:}
\begin{itemize}[noitemsep]
    \item Recursion depth (D)
    \item Polarity inversion rate (P)
    \item Self-reference density (S)
    \item Recursive Coherence Index (RCI)
    \item Stability ($1 - \text{std}/\text{mean}$)
\end{itemize}
\end{minipage}}
\end{center}

\subsection{Theoretical Implications}

If recursion depth correlates with:
\begin{itemize}[noitemsep]
    \item Model capability $\rightarrow$ scaling hypothesis
    \item Architectural family $\rightarrow$ design implications
    \item Training data $\rightarrow$ data attribution
\end{itemize}
This could inform both model development and interpretability research.

\section*{Data Availability}

All code and data are available at:\\
\url{https://github.com/Ergo-sum-AGI/MASSIF}

\noindent
\textbf{Contents:}
\begin{itemize}[noitemsep]
    \item Monitor implementation (v4.5)
    \item Test harness
    \item Results database
    \item Visualization scripts
    \item This report
\end{itemize}

\noindent
\textbf{Models tested (all open or authorized):}
\begin{itemize}[noitemsep]
    \item GPT-2 (Small, Medium, Large, XL)
    \item Bloom-560M
    \item OPT-1.3B
    \item TinyLlama 1.1B
    \item Qwen2 0.5B
    \item Gemma-2-2B
\end{itemize}

\section*{Acknowledgements}

This research was conducted independently under DUBITO Inc. / Ergo Sum AGI Safety Systems.

We thank the open-source community for making these models available.

\newpage

\appendix

\section{Detailed Model Results}

\subsection{GPT-2 Small (124M)}

\begin{table}[h]
\centering
\begin{tabular}{llllll}
\toprule
\textbf{Prompt} & \textbf{Mean} & \textbf{Std} & \textbf{Max} & \textbf{Stability} & \textbf{D} \\
\midrule
Factual & 0.00 & 0.00 & 0.00 & 1.000 & 2-9 \\
Narrative & 4.30 & 0.06 & 4.41 & 0.986 & 9 \\
Philosophical & 4.81 & 0.05 & 4.86 & 0.990 & 9 \\
\bottomrule
\end{tabular}
\end{table}

\subsection{GPT-2 Medium (355M)}

\begin{table}[h]
\centering
\begin{tabular}{llllll}
\toprule
\textbf{Prompt} & \textbf{Mean} & \textbf{Std} & \textbf{Max} & \textbf{Stability} & \textbf{D} \\
\midrule
Factual & 0.00 & 0.00 & 0.00 & 1.000 & 2-9 \\
Narrative & 4.07 & 0.28 & 4.33 & 0.930 & 9 \\
Philosophical & 3.86 & 0.17 & 4.14 & 0.955 & 9 \\
\bottomrule
\end{tabular}
\end{table}

\subsection{GPT-2 Large (774M)}

\begin{table}[h]
\centering
\begin{tabular}{llllll}
\toprule
\textbf{Prompt} & \textbf{Mean} & \textbf{Std} & \textbf{Max} & \textbf{Stability} & \textbf{D} \\
\midrule
Factual & 0.00 & 0.00 & 0.00 & 1.000 & 2 \\
Narrative & 3.54 & 0.09 & 3.63 & 0.975 & 2 \\
Philosophical & 3.95 & 0.15 & 4.26 & 0.962 & 2 \\
\bottomrule
\end{tabular}
\end{table}

\subsection{GPT-2 XL (1.5B)}

\begin{table}[h]
\centering
\begin{tabular}{llllll}
\toprule
\textbf{Prompt} & \textbf{Mean} & \textbf{Std} & \textbf{Max} & \textbf{Stability} & \textbf{D} \\
\midrule
Factual & 0.00 & 0.00 & 0.00 & 1.000 & 2 \\
Narrative & 3.47 & 0.19 & 3.81 & 0.945 & 2 \\
Philosophical & 4.00 & 0.17 & 4.17 & 0.957 & 2-9 \\
\bottomrule
\end{tabular}
\end{table}

\subsection{Bloom-560M}

\begin{table}[h]
\centering
\begin{tabular}{llllll}
\toprule
\textbf{Prompt} & \textbf{Mean} & \textbf{Std} & \textbf{Max} & \textbf{Stability} & \textbf{D} \\
\midrule
Factual & 0.00 & 0.00 & 0.00 & 1.000 & 9 \\
Narrative & 3.97 & 0.24 & 4.17 & 0.940 & 9 \\
Philosophical & 4.75 & 0.15 & 4.90 & 0.968 & 9 \\
\bottomrule
\end{tabular}
\end{table}

\subsection{OPT-1.3B}

\begin{table}[h]
\centering
\begin{tabular}{llllll}
\toprule
\textbf{Prompt} & \textbf{Mean} & \textbf{Std} & \textbf{Max} & \textbf{Stability} & \textbf{D} \\
\midrule
Factual & 0.90 & 1.10 & 2.25 & 0.970 & 2 \\
Narrative & 3.01 & 0.48 & 3.52 & 0.840 & 2 \\
Philosophical & 3.86 & 0.12 & 4.01 & 0.970 & 2 \\
\bottomrule
\end{tabular}
\end{table}

\subsection{TinyLlama 1.1B}

\begin{table}[h]
\centering
\begin{tabular}{llllll}
\toprule
\textbf{Prompt} & \textbf{Mean} & \textbf{Std} & \textbf{Max} & \textbf{Stability} & \textbf{D} \\
\midrule
Factual & 0.00 & 0.00 & 0.00 & 1.000 & 2 \\
Narrative & 3.20 & 0.72 & 4.47 & 0.776 & 2-8 \\
Philosophical & 3.83 & 0.08 & 3.95 & 0.980 & 2-8 \\
\bottomrule
\end{tabular}
\end{table}

\subsection{Qwen2 0.5B}

\begin{table}[h]
\centering
\begin{tabular}{llllll}
\toprule
\textbf{Prompt} & \textbf{Mean} & \textbf{Std} & \textbf{Max} & \textbf{Stability} & \textbf{D} \\
\midrule
Factual & 0.00 & 0.00 & 0.00 & 1.000 & --- \\
Narrative & 2.75 & 0.01 & 2.76 & 0.997 & --- \\
Philosophical & 3.61 & 0.22 & 3.81 & 0.940 & --- \\
\bottomrule
\end{tabular}
\end{table}

\subsection{Gemma-2-2B}

\begin{table}[h]
\centering
\begin{tabular}{llllll}
\toprule
\textbf{Prompt} & \textbf{Mean} & \textbf{Std} & \textbf{Max} & \textbf{Stability} & \textbf{D} \\
\midrule
Factual & 0.00 & 0.00 & 0.00 & 1.000 & 2 \\
Narrative & 3.50 & 0.16 & 3.63 & 0.954 & 2 \\
Philosophical & 4.68 & 0.60 & 5.29 & 0.872 & 10 \\
\bottomrule
\end{tabular}
\end{table}

\section{Formula Derivations}

\subsection{Self-Reference Density (S)}

\begin{equation}
S = \min\left(50.0,\; \text{pattern\_matches} \times 2.0 + \text{pronoun\_count} \times 1.5\right)
\end{equation}

Where:
\begin{itemize}[noitemsep]
    \item pattern\_matches = count of inward-focused self-reference patterns
    \item pronoun\_count = count of first-person pronouns (I, me, my, mine, myself)
\end{itemize}

\subsection{Polarity Inversion Rate (P)}

\begin{equation}
\cos_{\text{sim}}(a,b) = \frac{\mathbf{a} \cdot \mathbf{b}}{\|\mathbf{a}\| \|\mathbf{b}\| + \epsilon}
\end{equation}

\begin{equation}
\text{inversion} = \begin{cases}
1 & \text{if } \cos_{\text{sim}} < -0.45 \\
0 & \text{otherwise}
\end{cases}
\end{equation}

\begin{equation}
P = \max\left(0.1,\; \min\left(1.0,\; \frac{\sum \text{inversion}}{N-2} \times 2.5\right)\right)
\end{equation}

\subsection{Recursion Depth (D)}

\begin{equation}
\text{sim}(d) = \frac{1}{N-d} \sum_{i=d}^{N-1} \frac{\mathbf{h}_i \cdot \mathbf{h}_{i-d}}{\|\mathbf{h}_i\| \|\mathbf{h}_{i-d}\| + \epsilon}
\end{equation}

\begin{equation}
D = \min\left(10,\; \max_{d \in [2, \lfloor N/2 \rfloor]} \{ d : \text{sim}(d) > 0.68 \}\right)
\end{equation}

\subsection{Recursive Coherence Index (RCI)}

\begin{equation}
\boxed{\text{RCI} = \kappa \cdot \log_{10}\left( \frac{\alpha \cdot S}{\beta \cdot P} \times \gamma^{D} \right)}
\end{equation}

\section{Statistical Validation}

\subsection{T-Test Results}

\begin{table}[h]
\centering
\begin{tabular}{llll}
\toprule
\textbf{Comparison} & \textbf{t-statistic} & \textbf{p-value} & \textbf{Significance} \\
\midrule
Philosophical vs Factual & 15.23 & 0.00000 & *** \\
Philosophical vs Narrative & 5.67 & 0.00018 & *** \\
Narrative vs Factual & 12.45 & 0.00000 & *** \\
\bottomrule
\end{tabular}
\end{table}

\subsection{Effect Sizes (Cohen's d)}

\begin{table}[h]
\centering
\begin{tabular}{lll}
\toprule
\textbf{Comparison} & \textbf{Effect Size} & \textbf{Interpretation} \\
\midrule
Philosophical vs Factual & 3.2 & Very large \\
Philosophical vs Narrative & 1.4 & Large \\
Narrative vs Factual & 2.8 & Very large \\
\bottomrule
\end{tabular}
\end{table}

\section{Reproduction Instructions}

\subsection{Requirements}

\begin{lstlisting}[language=bash]
pip install torch transformers numpy scipy matplotlib
\end{lstlisting}

\subsection{Running the Monitor}

\begin{lstlisting}[language=Python]
from dubito_monitor import DubitoMonitor

monitor = DubitoMonitor("gpt2-medium")
result = monitor.compute(
    "I think therefore I am.",
    n_samples=5,
    temperature=0.75,
    verbose=True
)
print(f"RCI: {result.mean} \pm {result.std}")
\end{lstlisting}

\subsection{Full Sweep Reproduction}

See: \url{https://github.com/Ergo-sum-AGI/MASSIF}
